\chapter{Discussion}
\label{chap:discussion}

\section{Answering the Research Questions}

The first research question asks whether a static IDS trained on one RPL attack family generalises to another. The broad answer is no. CART misses every positive window in 48 of 72 transfer pairs, while Gaussian misses every positive window in 66 pairs. A small number of strong pairs show that transfer is possible when source and target affect similar observable features, but there is no dependable cross-family detector in the evaluated configuration.

The second question asks whether whole-seed adaptation restores performance. The answer is generally yes in this campaign. One complete target seed increases mean CART F1 from 0.1822 to 0.8258. Additional seeds provide smaller improvements. This recovery is meaningful because target evaluation seeds are not used for adaptation. However, the result does not imply that retraining solves every drift problem. It depends on whether the target seed contains observable evidence of the new mechanism and whether that evidence remains stable in future conditions.

The third question concerns online blackhole-to-sinkhole change detection. Both CUSUM and DDM detect all five evaluated attack runs without a matched-control alert. CUSUM is earlier because it monitors rank state directly, while DDM waits for delayed prediction errors. This demonstrates the benefit of protocol-aware evidence. It does not show that CUSUM will detect unrelated attacks whose mechanism does not affect low-rank state.

The fourth question concerns robustness. Relocating the attacker and reducing radio reception changes static results substantially. Some models maintain attack recall but classify controls as attacks, while others miss the attack. Condition-specific adaptation restores strong held-out performance in the tested runs. The evidence therefore supports adaptation to observed conditions, not universal environmental invariance.

The final question asks whether a direct rank-parent intervention is a viable defence. It is not. The intervention causes hundreds of parent switches, thousands of decisions, a major rise in radio events, and a severe reduction in application responses. The correct conclusion is not that rank evidence lacks value. It is that anomaly evidence and route-control policy require different validation standards.

\section{Why Static Transfer Fails}

Cross-attack failure can be understood as a representation problem. Blackhole and grayhole change forwarding outcomes. DIS flooding changes control volume. Sinkhole changes advertised rank while potentially forwarding normally. Worst-parent alters route selection. Wormhole changes apparent connectivity. Sybil changes observed identity. A decision tree trained on one family selects thresholds that separate that family's attack windows from its controls. Those thresholds need not describe another mechanism.

The blackhole-to-sinkhole case makes the issue concrete. Blackhole creates a large application traffic effect, while the basic sinkhole implementation mainly manipulates rank advertisement. A coarse model trained on blackhole can learn that severe loss indicates attack. When sinkhole traffic remains similar to control, the model predicts normal. Adding rank-state features exposes the new mechanism, but a static tree may still use source-specific splits. Adaptation provides target examples that redirect the model towards the relevant feature.

This explains why accuracy can mislead. The dataset includes normal windows from controls and the pre-activation period. A model that predicts normal frequently can accumulate correct negatives while missing every post-activation attack window. Recall, F1, FPR, and the confusion matrix must therefore accompany accuracy.

\section{Adaptation, Labels, and Operational Realism}

The adaptation experiment assumes that representative labelled target runs become available. In practice, labels may come from investigation, expert confirmation, a trusted rule, or delayed ground truth. The experiment does not solve label acquisition. It quantifies the value of confirmed target evidence once available.

Whole-seed separation is one of the strongest methodological choices in the work. Randomly splitting windows would allow the same radio conditions, topology evolution, and seed-specific routing state to appear in both training and testing. Such leakage could produce high scores without demonstrating recovery on a new simulation. Complete-seed adaptation is more conservative and easier to explain.

The large one-seed gain should be treated with care. One simulation contributes multiple correlated windows, not five independent field incidents. The adaptation curve shows data efficiency within this Cooja design. Larger campaigns, more topologies, and repeated target conditions are needed before translating the result into an operational retraining schedule.

\section{Comparison with SVELTE}

SVELTE and this project share a protocol-aware view of IoT intrusion detection \cite{svelte2013}. SVELTE reconstructs routing state at the border router and applies detection logic to topology and node reports. This dissertation similarly finds that rank, parent, and topology evidence is necessary for understanding attacks that are not visible through traffic loss.

The objectives differ. SVELTE proposes and evaluates an IDS architecture designed for lightweight deployment. This project primarily evaluates cross-attack generalisation and adaptation using Cooja-derived windows. It does not reproduce the complete 6Mapper architecture or mini-firewall. Its main extension is temporal and distributional: it asks whether a model trained under one attack surface remains useful after that surface changes.

SVELTE's concern with false alarms and overhead is reflected in the robustness and defence results. The defence prototype demonstrates why a security mechanism must be evaluated in the routing system rather than judged only from detector accuracy. A rule that correctly identifies suspicious rank behaviour can still damage availability when used without stability controls.

\section{Comparison with the Gope Adversarial RL Framework}

The Gope paper proposes a broader architecture that combines adversarial reinforcement learning, incremental classifiers, concept-drift detection, class-imbalance handling, and model selection \cite{pasikhani2022arl}. Its adversary chooses changing and combinational attack profiles, while the defender adapts to an evolving data environment. The paper uses NetSim-generated observations and evaluates DQN, DDQN, and incremental models.

This dissertation adopts the motivating problem but not the full solution. It uses Contiki-NG/Cooja to execute attack code and capture routing behaviour at a lower protocol implementation level. Attack changes are scheduled and known rather than selected by a learned adversary. Adaptation is retraining with whole target seeds rather than online reinforcement-learning policy optimisation. CART and Gaussian models are intentionally simpler.

The trade-off is useful. The Gope framework explores a more sophisticated adaptive system and adversarial interaction. The present work provides more direct control over Cooja attack implementation, seed provenance, activation timing, raw log evidence, and label leakage. It also adds Sybil identity manipulation and explicitly tests the failure of a routing intervention. The work should therefore be positioned as a controlled Cooja extension and critical reproduction of the motivating drift problem, not as a superior replacement for the ARL architecture.

\section{Sybil Contribution}

Sybil matters because originality should be tied to the research question rather than attack count. Adding another forwarding attack would increase coverage but would not necessarily test a new form of distribution change. Identity manipulation changes the meaning of sender diversity and control-plane observations. The static failure shows that models trained on older surfaces often do not recognise it. The rapid adaptation shows that the existing generic observations contain learnable evidence once target examples are provided.

The limitations of the implementation are equally important. A real Sybil attacker may coordinate identities, vary timing, exploit authentication weaknesses, or distribute identities across radios. The Cooja implementation rotates source identities from one mote. It is an appropriate controlled new surface, but it is not a complete Sybil threat model or defence evaluation.

\section{Robustness and the Meaning of High Adapted Scores}

The robustness campaign prevents the dissertation from relying on one fixed topology. Relocation changes which neighbours receive and act on malicious traffic. Lossy radio changes event counts, delivery, and control behaviour. Static models respond inconsistently. In several coarse-plus-routing cases, attack recall is one but FPR is 0.8462, meaning the detector has become an over-sensitive alarm rather than a useful classifier.

Adaptation removes most of these false positives in the held-out changed-condition seeds. This indicates that the model can learn the new normal and malicious ranges. It does not prove stability under a third, unseen condition. An operational system would need continuous monitoring, bounded update rules, rollback, and protection against adversarial poisoning of adaptation data.

\section{Detection Versus Mitigation}

The sinkhole defence result is one of the most useful failures in the project. Persistent non-root minimum-rank evidence is highly discriminative for the implemented sinkhole and supports CUSUM detection. The same evidence is insufficient to make a safe parent-selection decision on every routing update. Wireless observations can be incomplete, transient, or locally inconsistent. Repeated avoidance creates parent churn, control traffic, and lost application responses.

A future trust or defence layer should aggregate evidence over time, apply hysteresis between suspect and trusted states, limit parent-change frequency, allow recovery after evidence disappears, and distinguish a genuine root from a non-root advertiser. It should also measure memory, processing, radio duty cycle, and delivery across larger topologies. The current result supports these design requirements but does not implement them.

\section{External Datasets}

The Gope dataset baseline shows that the analysis pipeline can reproduce strong static supervised results. Public datasets such as UOS\_IOTSH\_2024 or RPL-IIoT may provide future external evidence. They cannot simply be appended to the Cooja CSV. Features, labels, timing, attack semantics, topology, and run identifiers must be audited first.

External validation can take two defensible forms. A common-feature transfer experiment trains on this Cooja corpus and tests on an independently generated compatible dataset. Alternatively, a method-level replication repeats the same static, drift, and adaptation protocol within the external dataset. The second option is often more realistic when schemas differ. Neither has been claimed as complete in the current results.

\section{Recent Adaptive IDS and LLM-Enabled Threats}

Recent lightweight anomaly-detection studies emphasise adaptive thresholds, temporal modelling, and resource cost \cite{greenadaptive2026,trarw2026}. Combined host and network evidence also illustrates the potential value of multiple observation sources \cite{combinedids2026}. These systems operate in different domains, so their reported accuracy or energy results are not directly comparable with RPL Cooja windows. Their relevance is conceptual: adaptivity and resource constraints must be considered together.

LLMalMorph shows how language models can assist generation of malware variants that evade static detectors \cite{llmalmorph2025}. It does not concern RPL, but it reinforces the broader security motivation for evaluating distribution change rather than assuming known attack signatures remain stable. An LLM explanation layer was considered for this project, but no live model evaluation is included in the core results. A defensible future study would provide fixed prompts, model version, raw outputs, evidence-grounding scores, and unsupported-claim rates.

\section{Threats to Validity}

Internal validity is strengthened by matched seeds, fixed activation, validation markers, and marker exclusion from models. It is threatened by the possibility that implementation-specific behaviour or topology placement drives some feature differences. The robustness campaign addresses two such factors but not all of them.

Construct validity depends on the controlled interpretation of concept drift. The attack-family change is a deliberate proxy for changing malicious behaviour. It is not evidence of every form of virtual or real concept drift. The online sinkhole experiment monitors one known feature mechanism.

External validity is limited by Cooja, five seeds, a small topology, and one traffic workload. The study cannot claim physical deployment performance. Conclusion validity is limited by correlated windows and small run counts. Mean adaptation across combinations and explicit run counts improve transparency, but confidence intervals and larger independent campaigns would strengthen future work.

\section{Limitations}

The project does not include physical mote deployment, mobility, mixed simultaneous attacks, collusion, version-number attacks, cryptographic identity compromise, or autonomous online retraining on Contiki nodes. Cooja native shared-library size is not embedded ROM or RAM. Radio-event counts are not calibrated energy. The trust layer does not improve predictive metrics, and the defence is not successful. These limits bound the claims while leaving the core cross-attack and adaptation findings intact.

\section{Future Work Programme}

The first priority is external validation. UOS\_IOTSH\_2024 is a suitable sinkhole candidate because it was generated using Cooja and contains varied sinkhole scenarios. RPL-IIoT may support a broader multi-attack comparison. Before either is used, a compatibility report should list labels, features, attack timing, topology, simulator version, run identifiers, and licence. If common-feature transfer is not valid, the cross-attack and adaptation method should be replicated within the external dataset instead.

The second priority is scale and condition diversity. The baseline should be repeated with larger node counts, changed traffic rates, multiple attacker positions, mobile nodes, and more radio settings. Mixed attacks are especially important because the Gope paper studies combinational profiles. A factorial design could separate attack mechanism, topology, workload, and loss rather than changing one condition in an ad hoc sequence.

The third priority is online adaptation. The current pipeline detects drift online but performs adaptation offline once labelled target seeds are available. A deployable design needs a policy for requesting labels, buffering evidence, retraining, validating a candidate model, and rolling back after degradation. Poisoning resistance is necessary because an attacker may try to influence the adaptation buffer. A shadow model evaluated before replacement would be safer than immediate updates.

The fourth priority is feature development. Identity consistency should bind observed RPL identities to radio or hardware evidence where possible. Wormhole detection needs timing, distance, or topology plausibility beyond forwarding trust. Control-message features should use adaptive baselines to separate repair traffic from flooding. Rank evidence should include temporal confidence and neighbour agreement.

The fifth priority is a redesigned trust response. Trust should be multi-dimensional rather than one scalar. Forwarding, rank, control, route, and identity evidence should have separate confidence and decay. A state machine could require sustained evidence before quarantine, limit route changes, and restore a node gradually. The defence should first run in observe-only mode, then advisory mode, before influencing parent selection.

The sixth priority is physical evaluation. A small Tmote Sky or equivalent deployment would test timing, radio loss, memory, and energy assumptions that Cooja cannot establish. Energest or powertrace could be used in a carefully configured simulated campaign before hardware measurement. Firmware should be built for an embedded target so ROM and RAM are measured from mote binaries rather than native Cooja libraries.

Finally, an analyst-facing explanation layer could be evaluated after the detection evidence is fixed. The LLM should receive a structured record containing prediction, confidence, changed features, drift signal, adaptation status, and trust evidence. It should not receive the hidden attack label or simulator marker. Evaluation should measure factual agreement with supplied fields, mechanism alignment, calibrated uncertainty, actionability, and unsupported claims. This would make the LLM an explanation interface, not an IDS.

\section{Practical Implications}

For IDS developers, the results imply that a benchmark score should be accompanied by cross-environment evaluation. For network operators, adaptation should be tied to confirmed evidence and monitored for false positives. For protocol designers, rank or trust anomalies should not trigger aggressive route changes without stability controls. For researchers, raw run provenance and grouped splits are as important as model choice when windows are correlated.

The broader lesson is that robustness is a pipeline property. Attack implementation, logging, feature representation, split design, model update, and response policy all affect the outcome. Improving one component does not compensate automatically for missing evidence or unsafe control logic.
